% p335_1; source PDF page 335; manually verified editable table.
\begin{tikzpicture}[x=1cm,y=-1cm,inner sep=0pt,outer sep=0pt]
\path[use as bounding box] (0,0) rectangle (16.000,20.825);
\draw[line width=.4pt] (0.000,0.000) rectangle (2.250,0.865);
\node[align=center,text width=1.990cm,font={\linespread{1}\fontsize{10}{13.50}\selectfont\bfseries}] at (1.125,0.432) {罪名};
\draw[line width=.4pt] (2.250,0.000) rectangle (4.500,0.865);
\node[align=center,text width=1.990cm,font={\linespread{1}\fontsize{10}{13.50}\selectfont\bfseries}] at (3.375,0.432) {成立条件};
\draw[line width=.4pt] (4.500,0.000) rectangle (11.650,0.865);
\node[align=center,text width=6.890cm,font={\linespread{1}\fontsize{10}{13.50}\selectfont\bfseries}] at (8.075,0.432) {行为对象和方式};
\draw[line width=.4pt] (11.650,0.000) rectangle (16.000,0.865);
\node[align=center,text width=4.090cm,font={\linespread{1}\fontsize{10}{13.50}\selectfont\bfseries}] at (13.825,0.432) {要点};
\draw[line width=.4pt] (0.000,0.865) rectangle (2.250,8.580);
\node[align=center,text width=1.990cm,font={\linespread{1}\fontsize{10}{13.50}\selectfont}] at (1.125,4.723) {生产、销售、提供假药罪};
\draw[line width=.4pt] (2.250,0.865) rectangle (4.500,8.580);
\node[align=center,text width=1.990cm,font={\linespread{1}\fontsize{10}{13.50}\selectfont}] at (3.375,4.723) {抽象危险犯（行为犯）};
\draw[line width=.4pt] (4.500,0.865) rectangle (11.650,8.580);
\node[align=left,text width=6.890cm,font={\linespread{1}\fontsize{10}{13.50}\selectfont}] at (8.075,4.723) {1.假药的判断标准，现采取实质标准，也即是否有疗效。有疗效，不算假药。以往采取形式标准，也即只要没有药监局批号，就是假药。\\{}2.假药的种类：\\{}（1）成分不符；\\{}（2）冒充的药品。以非药品冒充药品或者以他种药品冒充此种药品；\\{}（3）变质的药品；\\{}（4）标明的适应症超出规定范围。\\{}3.根据民间传统配方私自加工的药品，没有批号，但有疗效的，不算假药};
\draw[line width=.4pt] (11.650,0.865) rectangle (16.000,12.095);
\node[align=left,text width=4.090cm,font={\linespread{1}\fontsize{10}{13.50}\selectfont}] at (13.825,6.480) {1.提供假药罪、提供劣药罪，是指药品使用单位或其人员明知是假药、劣药而提供给他人使用。提供，包括有偿提供和无偿提供。有偿提供就是销售。\\{}2.不以营利为目的，实施带有自救、互助性质的生产、进口、销售药品的行为，不认定为犯罪。这里的药品包括假药、劣药及没有批号的药。\\{}3.未取得药品销售许可证，销售药品，以前定非法经营罪，现在不定非法经营罪。如果符合其他罪的要件，定其他罪};
\draw[line width=.4pt] (0.000,8.580) rectangle (2.250,12.095);
\node[align=center,text width=1.990cm,font={\linespread{1}\fontsize{10}{13.50}\selectfont}] at (1.125,10.338) {生产、销售、提供劣药罪};
\draw[line width=.4pt] (2.250,8.580) rectangle (4.500,12.095);
\node[align=center,text width=1.990cm,font={\linespread{1}\fontsize{10}{13.50}\selectfont}] at (3.375,10.338) {实害犯};
\draw[line width=.4pt] (4.500,8.580) rectangle (11.650,12.095);
\node[align=left,text width=6.890cm,font={\linespread{1}\fontsize{10}{13.50}\selectfont}] at (8.075,10.338) {劣药的种类：\\{}1.含量不符；\\{}2.被污染的药品；\\{}3.过期的药品；\\{}4.擅自添加防腐剂、辅料的药品};
\draw[line width=.4pt] (0.000,12.095) rectangle (2.250,18.510);
\node[align=center,text width=1.990cm,font={\linespread{1}\fontsize{10}{13.50}\selectfont}] at (1.125,15.303) {妨害药品管理罪};
\draw[line width=.4pt] (2.250,12.095) rectangle (4.500,18.510);
\node[align=center,text width=1.990cm,font={\linespread{1}\fontsize{10}{13.50}\selectfont}] at (3.375,15.303) {具体危险犯（足以严重危害人体健康）};
\draw[line width=.4pt] (4.500,12.095) rectangle (11.650,18.510);
\node[align=left,text width=6.890cm,font={\linespread{1}\fontsize{10}{13.50}\selectfont}] at (8.075,15.303) {1.生产、销售国务院药品监督管理部门禁止使用的药品的；\\{}2.未取得药品相关批准证明文件生产、进口药品或者明知是上述药品而销售的；（若适应症、成分不明，或在境外未合法上市，视为有具体危险）\\{}3.药品申请注册中提供虚假的证明、数据、资料、样品或者采取其他欺骗手段的；\\{}4.编造生产、检验记录的（若影响药品的安全性、有效性，视为有具体危险）};
\draw[line width=.4pt] (11.650,12.095) rectangle (16.000,18.510);
\node[align=left,text width=4.090cm,font={\linespread{1}\fontsize{10}{13.50}\selectfont}] at (13.825,15.303) {成立本罪，要求具体危险。\\{}私自进口、销售印度仿制药，如果没有具体危险，则不会构成该罪。因此，影片《我不是药神》中的男主角不会构成该罪，因为该药具有疗效};
\draw[line width=.4pt] (0.000,18.510) rectangle (2.250,20.825);
\node[align=center,text width=1.990cm,font={\linespread{1}\fontsize{10}{13.50}\selectfont}] at (1.125,19.668) {药品监管渎职罪};
\draw[line width=.4pt] (2.250,18.510) rectangle (4.500,20.825);
\node[align=center,text width=1.990cm,font={\linespread{1}\fontsize{10}{13.50}\selectfont}] at (3.375,19.668) {实害犯};
\draw[line width=.4pt] (4.500,18.510) rectangle (11.650,20.825);
\node[align=left,text width=6.890cm,font={\linespread{1}\fontsize{10}{13.50}\selectfont}] at (8.075,19.668) {本罪既包括故意（滥用职权），也包括过失（玩忽职守）};
\draw[line width=.4pt] (11.650,18.510) rectangle (16.000,20.825);
\node[align=left,text width=4.090cm,font={\linespread{1}\fontsize{10}{13.50}\selectfont}] at (13.825,19.668) {同时构成本罪和药品犯罪的共犯，想象竞合，择一重罪论处};
\end{tikzpicture}
